\documentclass[11pt]{letter}
\usepackage[utf8]{inputenc}
\usepackage[T1]{fontenc}
\usepackage{lmodern}
\usepackage{geometry}
\geometry{left=2.5cm,right=2.5cm,top=2.5cm,bottom=2.5cm}
\usepackage{amsmath,amsfonts}

\address{
Eduardo Di Santi \\
University of Colorado Boulder \\
Department of Computer Science \\
Boulder, CO 80309, USA \\
eduardo.disanti@colorado.edu
}

\signature{Eduardo Di Santi}

\begin{document}

\begin{letter}{}

\opening{Dear Editors,}

I am pleased to submit the manuscript entitled \textbf{``Cognitive Amplification vs Cognitive Delegation in Human--AI Systems: A Metric Framework''} for consideration as an Article in \textit{Nature Human Behaviour}.

This manuscript addresses a timely question in human--AI interaction: when does AI genuinely amplify human cognition, and when does it instead encourage progressive cognitive delegation? To address this problem, the paper introduces a compact operational framework built around four quantities---the Cognitive Amplification Index (\(CAI^*\)), the Dependency Ratio (\(D\)), the Human Reliance Index (\(HRI\)), and the Human Cognitive Drift Rate (\(HCDR\))---which together characterize collaborative gain, structural dependence, retained human contribution, and longitudinal cognitive change.

The paper combines conceptual analysis with an agent-based simulation in NetLogo. The simulation distinguishes several practically relevant regimes of human--AI interaction, including degenerate AI-dominated delegation, human-preserving but weakly competitive interaction, and intermediate boundary regimes that approach the AI baseline while remaining structurally dependent. A key result is that strong hybrid performance does not by itself imply genuine amplification: in the reported experiments, no tested regime achieved positive collaborative gain, and even a constrained optimization over capability atrophy failed to recover amplification under the most favorable simulated configuration. More broadly, the results suggest that genuine amplification may be considerably harder to obtain than conventional short-term performance measures would indicate.

I believe the manuscript is a strong fit for \textit{Nature Human Behaviour} because it addresses a broad interdisciplinary concern at the intersection of AI, cognition, decision-making, and human behaviour: how to evaluate AI systems not only by the outputs they help generate, but also by their effects on human capability over time. The framework is intended to be relevant across domains in which expertise retention, accountability, and cognitive sustainability matter, including medicine, engineering, and other safety-critical settings.

The manuscript is original, has not been submitted elsewhere, and I approve its submission. I would be happy to provide any additional material that may be useful for editorial assessment.

Thank you for your consideration.

\closing{Best regards,}

\end{letter}

\end{document}